%!TEX root = ../template.tex
\typeout{NT FILE chapter-workplan.tex}%


% ===================================================================
\chapter{Work Plan}
\label{ch:work-plan}
% ===================================================================
 
\section{Work Completed to Date}
\label{sec:wp-completed}
 
A fair share of this thesis's preparatory work happened before the
technology decisions in \S\ref{sec:pw-techstack} were even finalised,
so the phased plan below is laid out against what already exists.

The AR tracking evaluation reported in \S\ref{sec:pw-techstack} was not
desk research alone. All three candidate approaches, the hand-built
anchor-and-world-coordinate system, Vuforia, and Immersal, were actually
built and put through device tests, not just compared on paper. The Immersal
integration now runs end-to-end on a real Android device, localising
against a Chafariz Velho test map in roughly two to three seconds and
spawning test points of interest at their correct positions. This is
arguably the single most important piece of completed work, since it
takes this thesis's biggest technical risk off the table before
full-scale development even begins.

A second, fairly substantial body of completed work is the application's
visual identity, developed independently of any particular frontend
framework and so left untouched by the later move from Flutter to Unity.
The TileStories logo, colour palette, typography pairing (Fraunces for
editorial display text, DM Sans for interface text), and motion language
(entrance animations, press feedback, and a designed ``AR recognition''
moment where a gold ring pulses outward as the wall locks) were all
finalised during the Flutter prototyping phase as design decisions and carry forward unchanged into
the Unity build. The logo itself is shown alongside the application
concept in Figure~\ref{fig:logo} (\S\ref{subsec:pw-core-concept}).

A point-of-interest marker encoding system was also designed and
validated along the way, solving the problem of communicating a
building's category and its earthquake fate on a small map without the
two competing visually. Like the visual identity, this is a design
artefact that transfers to Unity directly; it does not need to be
redone. Finally, an earlier Flutter prototype and its accompanying
implementation guide were built before the tracking evaluation above was
even finished. Their specific code does not carry forward, but the UX and
domain-separation thinking behind them fed directly into
\S\ref{subsec:pw-core-concept} and \S\ref{sec:pw-architecture}. The
marker system, the visual identity, and screen-level prototypes from
this Flutter phase are presented together in Appendix~\ref{app:prototypes}.

\section{Development Phases Overview}
\label{sec:wp-overview}
 
Development is organised into the six phases introduced in
\S\ref{subsec:pw-feature-phases} (Table~\ref{tab:phases}): Core Tracking,
POI-Level Content, Scene-Level Features, Framework Packaging, Template
\& Third-Wall Validation, and Evaluation Readiness, plus a
time-permitting Stretch phase. The feature-level detail is not repeated
here; this section only adds the calendar dimension. Phase~0 also now
includes a task that is not yet resolved at the time of writing: a site
visit to capture reference photography of the Grande Panorama itself,
since the prior testing reported in \S\ref{sec:wp-completed} used
Chafariz Velho only. Whether this is even possible while MNAz remains
under renovation is still unresolved, not a scheduled certainty,
and is treated that way in the risk analysis below.
 
\section{Gantt Chart}
\label{sec:wp-gantt}
 
The six phases are laid out against a six-month calendar in
Appendix~\ref{app:gantt}, using relative month labels (M1--M6) instead
of fixed dates, since the confirmed start date depends on supervisor
sign-off following this report. Three points stand out even without
reproducing the chart itself. Phases overlap at
their edges on purpose, so evaluation findings from one phase can inform the next
before it fully begins. Phase~0 carries the MVP milestone, and Phase~4
is the point at which the framework is both packaged and validated on
the third wall. The stretch phase is scheduled, but only as
time-permitting; it is not a commitment.

\section{Risk Analysis}
\label{sec:wp-risks}

Two of the risks most likely to affect this thesis are access-related.
Both already have a working mitigation, not just a mention in passing:
not being able to photograph the Grande Panorama itself while MNAz
remains closed, and MNAz not reopening before submission at all. Neither
threatens the thesis's core claims, since Chafariz Velho and the Alto de
Santa Catarina mural (\S\ref{subsec:pw-chafariz},
\S\ref{subsec:pw-generalisability}) are both permanently and publicly
accessible, and together they are treated as a fully sufficient basis for
every stage of development and evaluation (\S\ref{sec:pw-problem-rq},
\S\ref{sec:pw-evaluation}). No contribution claim depends on Panorama
access. A third risk is more of a process bottleneck than a technical
one: the university's ethics approval for the street-recruited evaluation
(\S\ref{sec:pw-evaluation}) has not been requested yet, and, unlike the
museum-access risks, has no real fallback if it is slow, so it is
started as early in Phase~0 as possible rather than left until
Phase~5. A fuller risk table, also covering licensing cost, UI tooling
effort, recruitment shortfall, content-production pacing, and risks
specific to the framework's generalisability claim, is given in
Appendix~\ref{app:risks}.